% s14_t3_proof: the proof of prop:s14:T3 (input from sec/s14.tex)
\begin{proof}[Proof of \cref{prop:s14:T3}]
Units $m=1$, $N=2$, $M_1=0.7$; $\hat\mu=\lmu\in(0,1]$.

\emph{(i).} $S_{01,k}>0$ is \cref{lem:s14:A}(i). The four bounds with the interval upper sums are \cref{lem:s14:TF} with \cref{lem:s14:MON}, applied to the row of $\hat\mu$ (for $\hat\mu\ge10^{-4}$); the bounds with $B_3^*$ and $P^*+Q^*/\hat\mu$ are \cref{lem:s14:CF}; for $\hat\mu<10^{-4}$ the first bound is the $\mu$-free branch of $\bar j$: $j\le2M\bar\tau$ gives $\bar j_{400}(\hat\mu,\infty)\le\sum h\,(2M\bar\tau^*)^{\rm cell}=0.092610085$ for every $\hat\mu$, by the cell rule \eqref{eq:s14:cell}. The numbers $\bar j\le0.092610085<0.092611$ and $\bar s=6.5231964$ are read off \cref{tab:s14:T3}. The interval value $\bar B_4^-(\hat\mu)$, taken at $\mu_{\rm lo}$, and the closed form $P^*+Q^*/\hat\mu$, taken at $\hat\mu$, are two separate bounds, neither of which dominates the other on a whole row.

\emph{(ii).} By \eqref{eq:s14:struct} and \cref{lem:s14:A}(iii), \cref{lem:s14:B}(iii), \cref{lem:s14:C}(iii), the first four coefficients converge, so
\[
\sigma_1(0.7)=\liminf_ks^\infty_k(0.7)=A+b\hat\lambda-\hat e^2AI_\infty(\hat\mu)-L\hat\lambda^2+\liminf_k\big(\hat e^2\hat\lambda S_{21,k}+\hat e^4S_{40,k}\big),
\]
and by (i), for all but finitely many $k$ (those with $\eps_k\le1$), $\hat e^2\hat\lambda S_{21,k}+\hat e^4S_{40,k}\ge-\frac12\bar B_3(\hat\mu)\hat e^2\hat\lambda-\frac12\bar B_4^-(\hat\mu)\hat e^4$; with $I_\infty(\hat\mu)\le I_\infty(\mu_{\rm lo})$ (\cref{lem:s14:B}(iv)) this is $\sigma_1(0.7)\ge\sigma_{\rm low}$. The same argument with the closed-form bounds of (i) and $I_\infty(\hat\mu)\le I^*$ gives $\sigma_1(0.7)\ge\sigma^*$. (Cellwise, the interval value of $16\varphi_3$ at $\mu_{\rm lo}$ is at most the starred cell value, so $\bar B_3(\hat\mu)\le B_3^*$ always holds; for $b_4$ the interval value carries $\bar\delta(\mu_{\rm lo})\approx c_N/(8\pi\mu_{\rm lo})$ while the closed form carries $Q^*/\hat\mu$, so neither bound dominates the other on a whole row; hence the maximum in (ii).)

\emph{(iii).} The first sentence is (ii) and the definition of \Tthree\ (\cref{hyp:s03:T3}). $\sigma_{\rm low}$ is concave in $\hat\lambda$ at fixed $\hat e^2$ (the term $-L\hat\lambda^2$) and concave in $\hat e^2$ at fixed $\hat\lambda$ (the term $-\frac12\bar B_4^-\hat e^4$), so its minimum over a box $[0,c]^2$ is attained at a corner: the minimum over $\hat\lambda\in[0,c]$ of a concave function is attained at an endpoint, the two resulting functions of $\hat e^2$ are concave and so is their minimum, whose minimum over $[0,c]$ is again at an endpoint. The side $c(\hat\mu)$ of \cref{tab:s14:T3} is the largest $c$ for which the four corner values are $\ge0$, rounded down; $e^2_{\max}$ is the positive root of $\sigma_{\rm low}(0,\cdot)=A(1-\hat e^2I_\infty)-\frac12\bar B_4^-\hat e^4$, rounded down; and at $e=0$, $\sigma_{\rm low}=A+b\hat\lambda-L\hat\lambda^2>0$ if and only if $\hat\lambda<\lambda_c=\frac{b+\sqrt{b^2+4AL}}{2L}=2.14398\ldots$ (as in (vi)). \emph{Monotonicity in $\hat\mu$:} on $[10^{-4},1]$, $\sigma_{\rm low}(\cdot;\hat\mu)$ is a step function of $\hat\mu$, built row by row from $I_\infty(\mu_{\rm lo})$ and the interval values $\bar B_3$, $\bar B_4^-$ of the row, which enter with negative coefficients; so it suffices that these three row values do not decrease from each row to the row below it. For every row but the top one the interval values are the point values at $\mu_{\rm lo}$ (\cref{lem:s14:MON}, last clause, and the remark at the end of its proof for the row $[0.98,0.99]$), the point values are non-increasing in $\mu$ (\cref{lem:s14:MON}(i), with the point prefactor $\max_{M\in\{a,b\}}|\mu^2-2M^2|=2b^2-\mu^2$ non-increasing in $\mu$ for $\mu\le0.989949$) and so is $I_\infty$ (\cref{lem:s14:B}(iv)); hence the values of a row, taken at its $\mu_{\rm lo}$, are at most those of the row below, whose left end is smaller. For the top row $[0.99,1]$ this is read off the table ($8.5786210\le8.6036471$, $1.5854190\le1.6023646$, $0.008458\le0.008491$); by inspection each of the three columns is non-decreasing down the table through all $34$ rows. Across $\hat\mu=10^{-4}$: $\sigma^*$ is non-decreasing in $\hat\mu$, and $\sigma^*(\cdot;\hat\mu)\le\sigma_{\rm low}(\cdot;10^{-4})$ for $\hat\mu<10^{-4}$ because $I^*\ge I_\infty(10^{-4})$, $B_3^*\ge\bar B_3(10^{-4})=13.186961$ and $P^*+Q^*/\hat\mu>7716.84>\bar B_4^-(10^{-4})=7716.1679$. At a grid point both rows containing it give valid bounds, and the convention selects the row with $\mu_{\rm lo}=\hat\mu$. Hence $\hat\mu\mapsto\sigma_{\rm low}(\hat\lambda,\hat e^2;\hat\mu)$ is non-decreasing and $R(\lmu)\supseteq R(\lmu')$ for $\lmu\ge\lmu'$. \emph{The listed boxes:} $c\ge0.2063$ on the row $[0.5,0.55]$, $c\ge0.1525$ on $[0.1,0.15]$, $c\ge0.1132$ and $\sigma_{\rm low}(0.1,0.1)\ge0.0582$ on $[0.03,0.04]$, $c\ge0.0753$ on $[0.01,0.015]$ (\cref{tab:s14:T3}), together with the monotonicity; at $\hat\mu=1$ the row $[0.99,1]$ gives $c\ge0.2390$ and $e^2_{\max}\ge0.5618$.

\emph{(iv).} $\sigma^*$ is concave in each variable separately, so again the corners decide. Let $(\hat\lambda,\hat e^2)$ lie in the box $\{\hat\lambda\le\lambda_1,\ \hat e^2\le c_1\sqrt{\hat\mu}\}$ with $\hat\mu\le1$. At a corner $(\hat\lambda,c_1\sqrt{\hat\mu})$, $\hat\lambda\in\{0,\lambda_1\}$,
\[
\sigma^*=A-AI^*c_1\sqrt{\hat\mu}+b\hat\lambda-L\hat\lambda^2-\tfrac12B_3^*c_1\sqrt{\hat\mu}\,\hat\lambda-\tfrac12\big(P^*\hat\mu+Q^*\big)c_1^2\ \ge\ A-AI^*c_1+b\hat\lambda-L\hat\lambda^2-\tfrac12B_3^*c_1\hat\lambda-\tfrac12(P^*+Q^*)c_1^2=\sigma^*(\hat\lambda,c_1;1),
\]
using $\sqrt{\hat\mu}\le1$ and $\hat\mu\le1$ in the negative terms; the corners with $e=0$ give $A>0$ and $A+b\lambda_1-L\lambda_1^2>0$ ($\lambda_1\le1<\lambda_c$). So it suffices that $\sigma^*(0,c_1;1)\ge0$ and $\sigma^*(\lambda_1,c_1;1)\ge0$ for each listed pair, which is how the $c_1$ were determined (bisection on $c_1$, rounded down); explicitly, $\sigma^*(0,c_1;1)=A(1-I^*c_1)-\frac12(P^*+Q^*)c_1^2=3.8\cdot10^{-5}$ at $c_1=0.3014$, and $\sigma^*(\lambda_1,c_1;1)$ equals $1.2\cdot10^{-4}$, $4.1\cdot10^{-5}$, $1.1\cdot10^{-4}$, $2.3\cdot10^{-4}$, $6.3\cdot10^{-4}$ at the five pairs with $\lambda_1>0$, each positive.

\emph{(v).} \eqref{eq:s14:struct}, $\hat\lambda S_{01,k}\ge0$ (\cref{lem:s14:A}(i)) and (i).

\emph{(vi).} At $e=0$ only the orders $(0,1)$ and $(0,2)$ of \eqref{eq:s14:CTdef} are present, so by \cref{lem:s14:A}(iii), \cref{lem:s14:C}(iii) (for an arbitrary $M_1\in(0,m)$, and now in general units) $\lim_ks^\infty_k(M_1)=g_\lambda(M_1)\defeq(m-M_1)\big[\frac{m+M_1}2+\frac{(N+2)\lambda}{2\pi}\big]-\frac{(N+2)\lambda^2}{\pi^2}\ln\frac m{M_1}$ exactly, and $\sigma_1(M_1)=g_\lambda(M_1)$. As a function of $M_1\in(0,m]$: $g_\lambda(m)=0$, $g_\lambda'(M_1)=-M_1-\frac{(N+2)\lambda}{2\pi}+\frac{(N+2)\lambda^2}{\pi^2M_1}$, $g_\lambda''(M_1)=-1-\frac{(N+2)\lambda^2}{\pi^2M_1^2}<0$, so $g_\lambda$ is strictly concave. If $g_\lambda'(m)\ge0$, concavity gives $g_\lambda(M_1)<g_\lambda(m)+g'_\lambda(m)(M_1-m)\le0$ for every $M_1<m$, and \Tthree\ fails; if $g'_\lambda(m)<0$, then $g_\lambda>0$ on a left neighbourhood of $m$, and \Tthree\ holds. Now $g'_\lambda(m)<0$ if and only if $x=\lambda/m$ lies below the positive root of $\frac{N+2}{\pi^2}x^2-\frac{N+2}{2\pi}x-1$, which is $x_*=\frac\pi4\big(1+\sqrt{1+\frac{16}{N+2}}\big)$, equal to $\frac\pi4(1+\sqrt5)=2.5416\ldots$ for $N=2$. At a fixed $M_1$, $g_\lambda(M_1)>0$ if and only if $L\lambda^2-b\lambda-A<0$ with $A,b,L$ at $M_1$, i.e.\ $\lambda<\lambda_c(M_1)m\defeq\frac{b+\sqrt{b^2+4AL}}{2L}m$; at $M_1=0.7m$, $\lambda_c=2.14398\ldots$.

\emph{(vii).} By \cref{thm:a06:CD1}(b) the relative limit $\alpha_\infty(M)-\alpha_\infty(m)$ exists for every $M>0$ and $\sigma_1(M_1)=\lim_ks^\infty_k(M_1)=\frac12(\alpha_\infty(m)-\alpha_\infty(M_1))$. On each region of (iii), (iv), (vi), $\sigma_1(0.7m)>0$ (resp.\ $\sigma_1(M_1)>0$), so $\alpha_\infty(0.7m)<\alpha_\infty(m)$: $m$ is not a minimiser of $\alpha_\infty$ on $(0,m]$, which is \sxivNEl, and \sxivNEl\ implies \sxivNE\ (\cref{lem:s14:CW}(b)(i)). \Cref{prop:s10:Upm} is needed only on $E^<_W\setminus E_W$, which these regions avoid.
\end{proof}
